\section{Discussion and Design Implications}

\subsection{Coherence before precision}

Low precision changes an iterative solver in two ways. It changes arithmetic
values, and it can change the decisions made from those values. In pair-SMO,
a small perturbation can select a different pair; that pair changes the next
gradient, which changes every subsequent selection. Precision is consequently
not a local property of the GEMM that produces a dot product or of the cache
that stores an RBF value. It is a property of the induced state-transition
system.

This explains why converting only the cache can be both smaller and slower.
If the local solver selects pairs using one representation while the global
gradient is updated with another, the implementation repeatedly reconciles
two different operators and may follow a less productive trajectory. Coherent
FP16 succeeds by making the materializer, cache, and decision-making consumers
agree on the same rounded kernel. It does not require every scalar control
operation to use FP16. The transferable rule is therefore
\emph{coherence before precision}: first define which approximate operator the
algorithm solves, then choose the precision of each operation while preserving
that operator at every state-changing boundary.

\subsection{Progress per kernel-data item}

The conventional metrics for a GPU kernel---throughput, occupancy, bytes per
second, or pair updates per second---do not measure whether an optimization
solver is moving efficiently toward its target. In this workload, a global
column $K[X,j]$ may require a full $N\times d$ feature scan, whereas a local
row $K[W,j]$ is read from an already materialized working-set matrix. Counting
both as one ``kernel row'' obscures the dominant cost.

A more useful comparison couples the weighted work in
\cref{eq:cost-model} with optimization progress. For example, over an interval
$[t,t+r]$, one can report
\begin{equation}
  \eta_{\mathrm{kernel}}
  =
  \frac{\Phi(\bm{\alpha}^{t})-\Phi(\bm{\alpha}^{t+r})}
       {M_g C_g(N,d)+R_\ell C_\ell(|W|)},
  \label{eq:progress-per-kernel-data}
\end{equation}
where $\Phi$ is the minimized dual objective and the denominator weights
global-column generation separately from local-row use. KKT-gap reduction can
be substituted when objective values are unavailable. This metric exposes a
common failure mode: a speculative method can execute more GPU work per
second while obtaining less progress from every generated column.

The resulting design goal is not simply to maximize cache hits or reduce pair
latency. It is to reuse every expensive, immutable kernel column for enough
useful local progress to amortize its generation, without over-solving a stale
working set. This framing applies more broadly to iterative GPU algorithms in
which expensive data products feed a sequence of state-dependent decisions.

\subsection{Work--depth, not occupancy}

The serial chain in \cref{eq:serial-chain} creates low device utilization, but
launching more work does not remove the dependence. Large-block methods,
fixed-depth replay, and speculative prefetch can expose more parallelism while
also introducing larger Hessians, additional local rows, wasted global
columns, projection work, or extra outer rounds. Occupancy improves even when
time to convergence does not.

The relevant abstraction is a work--depth frontier. Here, work includes the
first two terms of \cref{eq:cost-model}, while depth counts state-changing
decision stages that cannot overlap. A useful parallel solver must reduce
depth faster than it increases weighted kernel-data work. The fixed-depth and
triangular experiments fail this test: they shorten or regularize control but
leave too much row work unchanged. This is a stronger conclusion than
``the GPU is underutilized.'' It identifies what a replacement algorithm must
change: the optimization directions themselves must yield more progress per
dependent stage.

\subsection{Shape and phase are first-class inputs}

The profitable execution regime is nonuniform in both data shape and solver
state. At $d=18$ or 28, padding, repeated local-row loads, and solver control
can dominate because each kernel evaluation contains little arithmetic. At
$d=2{,}381$, Tensor Core dot products, full-data column misses, cache capacity,
and the local solver are all material. A representation or tiling strategy
that is effective in one regime can add overhead in the other.

The same nonuniformity appears over time. In the measured exact trace, late
outer rounds 100--140 contain residual-limited runs with a local
5$\times$--7$\times$ depth opportunity, whereas early and middle rounds are
frequently interrupted by box constraints. This suggests phase-aware solver
selection, but it does not make adaptation free: each phase-specific primitive
must first be profitable on its own, and the switching signal must cost less
than the work it avoids.

Taken together, these results favor dispatch over a universal kernel. Feature
dimension, working-set size, cache state, active constraints, and recent
progress should select among mechanisms only when each mechanism has a
measured region of advantage.

\subsection{Amdahl limits and co-design priorities}

The steady solver phase in \cref{tab:profile} makes the remaining limits
concrete. Completely deleting the 44.55 ms materializer would provide at most
approximately 1.27$\times$ speedup. Completely deleting the 84.75 ms block
solver would provide at most approximately 1.68$\times$. A realistic
2$\times$ improvement in only one stage yields approximately 1.12$\times$ and
1.25$\times$, respectively. Thus, no remaining microkernel alone can deliver
a 2$\times$ end-to-end gain on this workload.

The more promising target is an algorithm--system co-design that changes both
terms: use a representation that increases reuse and lowers column cost, then
use a solver that obtains more progress before the next full-data update. Such
a design should be judged by three questions: how many new global columns it
generates, how many local rows it consumes, and how much objective or KKT
progress it obtains before the next dependent decision. These questions turn
the present negative results into reusable gates for future designs rather
than implementation-specific anecdotes.
